\documentclass[a4paper, 12pt]{article}
\usepackage[english, russian]{babel}
\usepackage[T2A]{fontenc}
\usepackage[utf8]{inputenc}
\usepackage{indentfirst}
\usepackage{geometry}
\usepackage{titleps}
\usepackage{mathtext}
\usepackage{amsmath}
\usepackage{enumitem}
\usepackage{tikz}
\usepackage{needspace}
\setlength{\tabcolsep}{4pt}
\newpagestyle{main}{
\setheadrule{0pt}
\sethead{}{}{}
\setfootrule{0pt}
\setfoot{}{}{}
}
\pagestyle{main}

\usepackage{amsthm}
\geometry{top=2cm}

% Блоки схем алгоритмов рисуются в масштабе 1 клетка = 0.5 см
\tikzset{blk/.style={x=0.5cm, y=0.5cm, line width=0.6pt}}

\begin{document}
\begin{center}
{\large КОНСПЕКТ ЛЕКЦИЙ}

ПО ДИСЦИПЛИНЕ <<ПРОГРАММИРОВАНИЕ>>
\end{center}
\vspace{10pt}

\section*{Арифметические выражения}

Операция: \texttt{17 + 2}. Оператор (команда, инструкция).

Арифметическое выражение может состоять из:
\begin{enumerate}[noitemsep, label=\arabic*)]
  \item операций:
    \begin{center}
    \begin{tabular}{cl}
      \texttt{+}  & сложение \\
      \texttt{-}  & вычитание \\
      \texttt{*}  & умножение \\
      \texttt{/}  & деление \\
      \texttt{//} & целочисленное деление \\
      \texttt{\%} & остаток от деления \\
      \texttt{**} & возведение в степень \\
    \end{tabular}
    \end{center}
  \item скобок:
    \[
      \frac{a + b}{c + d} \;\Longrightarrow\; \texttt{(a + b) / (c + d)}
    \]
  \item переменных;
  \item числовых констант;
  \item математических функций:
    \begin{center}
    \begin{tabular}{ll}
      1 способ & 2 способ \\
      \texttt{import math} & \texttt{from math import *} \\
      \texttt{math.log(4.8)} & \texttt{log(4.8)} \\
    \end{tabular}
    \end{center}
  \item $\sqrt{x}$ --- \texttt{sqrt(x)}; \quad \texttt{x ** (1/2)} --- комплексное число!
    (если введём отрицательное число);
  \item тригонометрических функций --- аргумент в радианах;
  \item \texttt{degrees} --- перевод из радиан в градусы;
    \texttt{radians} --- перевод $x$ из градусов в радианы.
\end{enumerate}

\subsection*{Приоритет операций}

\begin{enumerate}[noitemsep, label=\arabic*)]
  \item то, что в \texttt{( )};
  \item математические функции;
  \item[3--4)] операция возведения в степень; унарная операция \texttt{(-)}:
    \texttt{-5}, \texttt{-sin(x)}, \texttt{+5} (в мат. --- комплементарная); % так в оригинале
  \item[5)] \texttt{*} \ \texttt{/} \ \texttt{\%} \ \texttt{//};
  \item[6)] \texttt{+} \ \texttt{-}.
\end{enumerate}

\texttt{x = 1 + 2 + 3} --- по очереди слева направо.

\texttt{2 ** 3 ** 2} $\to$ \texttt{512} --- справа налево.

\texttt{\#} --- однострочный комментарий.

\texttt{'''} многострочный комментарий \texttt{'''}.

\section*{Тема 2. Линейные алгоритмы}

\textbf{Алгоритм} --- это конечная последовательность действий решения
какой-либо задачи.

Способы представления алгоритма:
\begin{enumerate}[noitemsep]
  \item структурограммы;
  \item схема алгоритма;
  \item псевдокод;
  \item словесная запись.
\end{enumerate}

\subsection*{Основные блоки схемы алгоритма}

\begin{enumerate}[label=\arabic*)]
  \item \Needspace{3cm}Вычислительный процесс
    \begin{center}
    \begin{tikzpicture}[blk]
      \draw (0,0) rectangle (6,2);
      \node[above] at (3,2) {6 кл};
      \node[left] at (0,1) {2};
    \end{tikzpicture}
    \end{center}

  \item \Needspace{3cm}Начало и конец процесса
    \begin{center}
    \begin{tikzpicture}[blk]
      \draw[rounded corners=0.5cm] (0,0) rectangle (6,2);
    \end{tikzpicture}
    \end{center}

  \item \Needspace{3cm}Ввод и вывод
    \begin{center}
    \begin{tikzpicture}[blk]
      \draw (1,0) -- (7,0) -- (6,2) -- (0,2) -- cycle;
    \end{tikzpicture}
    \end{center}

  \item \Needspace{3cm}Решение
    \begin{center}
    \begin{tikzpicture}[blk]
      \draw (0,1) -- (3,2) -- (6,1) -- (3,0) -- cycle;
      \node[above] at (3,2) {6};
      \node[left] at (0,1) {2};
    \end{tikzpicture}
    \end{center}

  \item \Needspace{3cm}Цикл с параметром (старый ГОСТ)
    \begin{center}
    \begin{tikzpicture}[blk]
      \draw (0,1) -- (1,2) -- (5,2) -- (6,1) -- (5,0) -- (1,0) -- cycle;
      \node[above] at (3,2) {6};
      \node[left] at (0,1) {2};
    \end{tikzpicture}
    \end{center}

  \item \Needspace{3cm}Цикл с параметром (новый ГОСТ)
    \begin{center}
    \begin{tikzpicture}[blk]
      \draw (0,0) -- (6,0) -- (6,1.5) -- (5.5,2) -- (0.5,2) -- (0,1.5) -- cycle;
    \end{tikzpicture}
    \qquad
    \begin{tikzpicture}[blk]
      \draw (0,0.5) -- (0.5,0) -- (5.5,0) -- (6,0.5) -- (6,2) -- (0,2) -- cycle;
    \end{tikzpicture}
    \end{center}

  \item \Needspace{3cm}Соединитель
    \begin{center}
    \begin{tikzpicture}[blk]
      \draw (1,1) circle (1);
      \node[above] at (1,2) {2};
      \node[left] at (0,1) {2};
    \end{tikzpicture}
    \end{center}

  \item \Needspace{3cm}Межстраничный соединитель
    \begin{center}
    \begin{tikzpicture}[x=1cm, y=1cm, line width=0.6pt]
      \draw (0,1) -- (1,1) -- (1,0.35) -- (0.5,0) -- (0,0.35) -- cycle;
      \node[above] at (0.5,1) {1};
      \node[left] at (0,0.5) {1};
    \end{tikzpicture}
    \end{center}

  \item \Needspace{3cm}Предопределённый процесс
    \begin{center}
    \begin{tikzpicture}[blk]
      \draw (0,0) rectangle (6,2);
      \draw (0.5,0) -- (0.5,2) (5.5,0) -- (5.5,2);
      \node[above] at (3,2) {6};
      \node[left] at (0,1) {2};
    \end{tikzpicture}
    \end{center}
\end{enumerate}

\subsection*{Правила составления}

\begin{enumerate}[noitemsep, label=\arabic*)]
  \item Одинаковая ширина блоков --- 6--8 клеток.
  \item Высота блоков может быть различной.
  \item Соединительная стрелка между блоками --- 1 клетка, по необходимости 2.
  \item Математические функции пишем как в математике.
  \item Подряд идущие ввод данных, вывод данных, присваивание объединяем
    в один блок.
  \item Каждое присваивание в блоке пишем с новой строки.
  \item Схема алгоритма изображается без привязки к конкретному языку
    программирования.
  \item Схема алгоритма должна занимать в высоту не более листа А4.
\end{enumerate}

\end{document}
